\section{Results}
\label{sec:results}

This section reports the results of the study. In \cref{subsec:training_results} we assess the accuracy of the trained LDNet on the held-out test set and document the hyperparameter selection. In \cref{subsec:mpc_results} we evaluate the closed-loop controller over the CS-25 discrete-gust envelope in nominal conditions, and in \cref{subsec:noise_results,subsec:systematic_results} we probe its robustness to stochastic measurement noise and to systematic errors of the gust preview.

\subsection{LDNet accuracy and hyperparameter tuning}
\label{subsec:training_results}

We swept the degrees of freedom that dominate the validation loss --- the latent dimension $d_s$, the depth of $\mathcal{NN}_\mathrm{dyn}$ and $\mathcal{NN}_\mathrm{rec}$, and the composition of the input vector $\boldsymbol{\eta}$ --- while holding the remaining hyperparameters fixed. Along each swept axis the candidate configurations are trained in multiple independent runs that differ only in the random initialisation of the weights, every run is scored by the loss on the validation set, and the run with the lowest validation loss is retained, so that the comparison is guarded against the sensitivity of the full-batch optimisation to its trajectory.

Following~\cite{Regazzoni}, accuracy is measured by the normalised root-mean-square error and by the Pearson dissimilarity $1-\rho$. For a scalar load $y \in \{F_y, M_z\}$, a component of the output $\mathbf{y}$, with prediction $\tilde{y}$ over the $N$ test samples and time steps,
\begin{equation}
\mathrm{NRMSE}(y) =
\frac{
\sqrt{\frac{1}{N}\sum_k \left(\tilde{y}_k - y_k\right)^2}
}{
\max_k y_k - \min_k y_k
},
\qquad
\rho(y) =
\frac{
\mathrm{cov}(\tilde{y}, y)
}{
\sigma_{\tilde{y}}\, \sigma_y
},
\label{eq:acc_metrics}
\end{equation}
so that both $\mathrm{NRMSE}$ and $1-\rho$ vanish for a perfect reconstruction. The two loads are also reported jointly as a combined value, obtained by stacking $F_y$ and $M_z$ and normalising by a single global range. Because the loads carry different units (N and N\,m), this stacking is not a physically meaningful norm and the combined value tracks the wider-ranged $F_y$: within a sweep we use it only as a single-number ordering criterion, always alongside the per-load values in the same table, and elsewhere we quote it only as a compact summary figure for the retained configuration, never as a substitute for the per-load NRMSE when what is being reported is an accuracy level to be interpreted on its own. The combined $1-\rho$ is a more severe case of the same stacking artefact --- normalising by a single global range collapses it roughly an order of magnitude below either per-load value --- so we do not quote it as an accuracy figure at all, only $\mathrm{NRMSE}$.

We deliberately perform the analysis of this subsection in open loop: the networks are driven by the reference structural signals of each trajectory, in the teacher-forcing regime of \cref{subsec:training}, and we measure the reconstruction error of the loads on the held-out test set by the metrics~\eqref{eq:acc_metrics}. Decoupled from the error propagation of an autonomous rollout, this error isolates how accurately the LDNet has learned the load mapping, and it is on this learning accuracy that the selection of the inputs and of $d_s$ is based. The regularisation weights $\alpha_\mathrm{dyn}=\alpha_\mathrm{rec}=\alpha_\mathrm{reg}=3\cdot10^{-4}$ are held equal and fixed at this value throughout the sweeps below, and Tikhonov regularisation of the network weights~\cite{Regazzoni} adds to the training loss a penalty term, the mean of the squared kernels of $\mathcal{NN}_\mathrm{dyn}$ and $\mathcal{NN}_\mathrm{rec}$ weighted by $\alpha_\mathrm{reg}$. Training stops once the validation loss, computed without that term, has not improved for $200$ iterations, and the weights retained are those that attained its minimum. The closed-loop rollout is instead the mode in which the LDNet is used: trained on the rollout loss~\eqref{eq:rollout_loss} with the damped latent update~\eqref{eq:latent_dyn}, the LDNet runs in rollout inside the control loop of the GLA study of \cref{subsec:mpc_results}.

\subsubsection{Input selection (Study 1)}
\label{subsubsec:input_selection}

Study~1 fixes the outputs $(F_y,M_z)$ and varies the signals fed to the networks, to
establish how much of the structural state the load reconstruction requires. Three
configurations are compared: gust and flap only $\{W,\delta\}$; the same augmented with the
plunge and pitch $\{W,\delta,h,\alpha\}$; and the full structural state
$\{W,\delta,h,\dot h,\alpha,\dot\alpha\}$, which adds the plunge and pitch rates.
\Cref{fig:input_study} reports NRMSE and $1-\rho$ against the latent dimension $d_s$ for
the three sets, aggregated over the held-out test trajectories of the three dataset families
of \cref{tab:dataset_families}, which span gust-only, flap-only and combined gust\,+\,flap
excitation. The gust-and-flap pair alone is inadequate: without the structural state the
latent dynamics must reproduce the aeroelastic response unaided, and the combined NRMSE is the
largest of the three sets.
Adding the plunge and pitch positions stabilises the reconstruction, but the four-signal set
does not improve beyond $4.0\cdot10^{-3}$ at any latent dimension. The velocity rates
$\dot h,\dot\alpha$ are what admit the accurate regime: only the full six-signal set reaches a
combined NRMSE of $1.3$--$1.4\cdot10^{-3}$. We
adopt the full six-signal set for the remainder of this work.

\begin{figure}[t]
\centering
\includegraphics[width=\textwidth]{Thesis_article_format/Chapters/Chapter2/Images/fig_ch2_input_study.png}
\caption{Study~1: open-loop test-set accuracy of the load reconstruction (NRMSE, left; Pearson
dissimilarity $1-\rho$, right; combined value) versus latent dimension
$d_s\in\{1,3,5,10\}$, for the three input configurations $\{W,\delta\}$,
$\{W,\delta,h,\alpha\}$ and the full structural state
$\{W,\delta,h,\dot h,\alpha,\dot\alpha\}$. Each point aggregates the held-out test
trajectories of the three dataset families of \cref{tab:dataset_families}. Only the full
six-signal set reaches the accurate regime, a combined NRMSE of $1.3$--$1.4\cdot10^{-3}$.}
\label{fig:input_study}
\end{figure}

\subsubsection{Latent-state dimension (Study 2)}
\label{subsubsec:latent_dim}

With the six-signal input fixed, Study~2 sweeps the latent dimension over $d_s\in\{1,3,5,10\}$; for
each value two independent training runs are performed and the run with the
lowest validation loss is retained. \Cref{fig:input_study} of Study~1 and \cref{tab:latent_sweep}
report the resulting test-set accuracy. A single latent state caps the combined NRMSE at
$4.6\cdot10^{-3}$; with more states the optimiser can reach a deep minimum of the validation
loss ($\mathcal{L}_\mathrm{val}\approx 2\cdot10^{-5}$) where the combined NRMSE drops to $1.4\cdot10^{-3}$ at $d_s=3$
and $1.3\cdot10^{-3}$ at $d_s=10$, with $1-\rho\le2.8\cdot10^{-4}$ on $F_y$ and $\le2.0\cdot10^{-4}$ on $M_z$
at both. The deep minimum is
not reached in every run --- at $d_s=5$ neither of its two runs located it --- a residue of the
sensitivity of the full-batch optimisation to its trajectory, which the multiple-run protocol
mitigates. We set the latent dimension to $d_s=10$, the most accurate configuration, which is
also the value adopted for the closed-loop rollout (\cref{subsec:mpc_results}).

Latent dimensions above $d_s=10$, up to $d_s=20$, were also probed without further gain, and the
reason is architectural rather than statistical: with two hidden layers of seven neurons,
$\mathcal{NN}_\mathrm{dyn}$ bounds the Jacobian of the latent recurrence~\eqref{eq:latent_dyn} to
rank seven, so for $d_s>7$ at least $d_s-7$ latent directions cannot be driven by the network at
all --- already three of the ten retained above. Left undriven, their eigenvalue in the recurrence
sits at $1-\lambda$ --- unity with the
decay switched off --- rather than at a value training can shape, so widening $d_s$ past that
point adds inert state, not usable dynamics. Reaching further requires widening
$\mathcal{NN}_\mathrm{dyn}$ itself; \cref{subsubsec:depth} (Study~3) instead sweeps its depth at
fixed width.


\begin{table}[t]
\centering
\begin{tabular}{|c|c|c|c|c|c|c|}
\hhline{|-|-|-|-|-|-|-|}
\rowcolor{bluePoli!40}
\multicolumn{1}{|c|}{} & \multicolumn{3}{c|}{NRMSE} & \multicolumn{3}{c|}{$1-\rho$} \T\B \\
\hhline{|~|-|-|-|-|-|-|}
\rowcolor{bluePoli!40}
\multicolumn{1}{|c|}{$d_s$} & \multicolumn{1}{c|}{$F_y$} & \multicolumn{1}{c|}{$M_z$} & \multicolumn{1}{c|}{comb.} & \multicolumn{1}{c|}{$F_y$} & \multicolumn{1}{c|}{$M_z$} & \multicolumn{1}{c|}{comb.} \T\B \\
\hhline{|=|=|=|=|=|=|=|}
$1$  & $8.8\cdot10^{-3}$ & $2.9\cdot10^{-3}$ & $4.6\cdot10^{-3}$ & $2.6\cdot10^{-3}$ & $6.2\cdot10^{-4}$ & $2.0\cdot10^{-4}$ \T\B \\
\hline
$3$  & $2.8\cdot10^{-3}$ & $1.6\cdot10^{-3}$ & $1.4\cdot10^{-3}$ & $2.8\cdot10^{-4}$ & $2.0\cdot10^{-4}$ & $2.0\cdot10^{-5}$ \T\B \\
\hline
$5$  & $9.3\cdot10^{-3}$ & $4.8\cdot10^{-3}$ & $4.8\cdot10^{-3}$ & $3.2\cdot10^{-3}$ & $1.8\cdot10^{-3}$ & $2.2\cdot10^{-4}$ \T\B \\
\hline
$10$ & $2.4\cdot10^{-3}$ & $1.6\cdot10^{-3}$ & $1.3\cdot10^{-3}$ & $2.1\cdot10^{-4}$ & $2.0\cdot10^{-4}$ & $1.5\cdot10^{-5}$ \B \\
\hline
\end{tabular}
\\[6pt]
\caption{Study~2: open-loop test-set reconstruction accuracy of the aerodynamic loads as a function of
the latent dimension $d_s$, with the six-signal input: NRMSE and Pearson dissimilarity $1-\rho$ for
$F_y$, $M_z$ and the combined output. For each $d_s$ the best of the independent training runs,
selected on the validation loss, is reported.}
\label{tab:latent_sweep}
\end{table}

\subsubsection{Network depth (Study 3)}
\label{subsubsec:depth}

Study~3 addresses the last architectural degree of freedom, the depth of the two networks. The layer widths are
held fixed at $7$ neurons for $\mathcal{NN}_\mathrm{dyn}$ and $24$ for $\mathcal{NN}_\mathrm{rec}$,
and the total number of hidden layers $L$ is swept starting from $L=6$, the depth used elsewhere
in this work, doubled twice over $L\in\{6,12,24\}$ in the fixed $1{:}2$ split between the two
networks; \cref{tab:depth_arch} lists the three configurations. Every depth is trained at the two
latent dimensions that bracket the sweep of \cref{tab:latent_sweep}, $d_s\in\{1,10\}$, each
configuration twice from independent weight initialisations so that the effect of depth can be
told apart from the run-to-run scatter of the latent sweep, with the six-signal input, the
regularised trainer and the metrics of \cref{eq:acc_metrics} used throughout; \cref{tab:depth_loads}
collects the test-set error of the twelve runs.

Doubling the depth from $L=6$ improves the open-loop reconstruction at both latent dimensions:
$L=12$ reaches the most accurate configuration of the entire hyperparameter study, and at
$d_s=10$, for the two initialisations run at each configuration, both reached that accuracy
level while at $L=6$ only one of the two did; at $d_s=1$ the two $L=6$ seeds are already close
($3.66\cdot10^{-3}$ and $3.68\cdot10^{-3}$). Depth does not, however, replace latent capacity — at $d_s=1$ the
deeper network still trails the shallow one at $d_s=10$ — so the latent dimension sets the
accuracy level, and, within the two seeds tested per configuration, depth is what let training
reach it consistently. The depth sweep stops at
$L=24$: a plain $\tanh$ stack of this depth, without residual connections or a depth-aware
initialisation, saturates and attenuates the gradient layer after layer~\cite{glorot2010}, and
neither seed reaches a usable minimum, undoing the gains seen at $L=12$.

\begin{table}[t]
\centering
\begin{tabular}{|c|c|c|c|c|}
\hhline{|-|-|-|-|-|}
\rowcolor{bluePoli!40}
\multicolumn{1}{|c|}{} & \multicolumn{2}{c|}{hidden layers $\times$ width} & \multicolumn{2}{c|}{$P$} \T\B \\
\hhline{|~|-|-|-|-|}
\rowcolor{bluePoli!40}
\multicolumn{1}{|c|}{$L$} & \multicolumn{1}{c|}{$\mathcal{NN}_\mathrm{dyn}$} & \multicolumn{1}{c|}{$\mathcal{NN}_\mathrm{rec}$} & \multicolumn{1}{c|}{$d_s=1$} & \multicolumn{1}{c|}{$d_s=10$} \T\B \\
\hhline{|=|=|=|=|=|}
$6$  & $2\times7$ & $4\times24$  & $2217$ & $2568$  \T\B \\
\hline
$12$ & $4\times7$ & $8\times24$  & $4729$ & $5080$  \T\B \\
\hline
$24$ & $8\times7$ & $16\times24$ & $9753$ & $10104$ \B \\
\hline
\end{tabular}
\\[6pt]
\caption{Study~3: architectures of the depth sweep: total number of hidden layers $L$, number and width
of the hidden layers of each network, and number of trainable weights $P$ at the two latent
dimensions swept. $L=6$ is the depth used elsewhere in this work, and the sweep's starting point.}
\label{tab:depth_arch}
\end{table}

\begin{table}[t]
\centering
\begin{tabular}{|c|c|c|c|c|}
\hhline{|-|-|-|-|-|}
\rowcolor{bluePoli!40}
\multicolumn{1}{|c|}{} & \multicolumn{2}{c|}{$d_s=1$} & \multicolumn{2}{c|}{$d_s=10$} \T\B \\
\hhline{|~|-|-|-|-|}
\rowcolor{bluePoli!40}
\multicolumn{1}{|c|}{$L$} & \multicolumn{1}{c|}{best} & \multicolumn{1}{c|}{worst} & \multicolumn{1}{c|}{best} & \multicolumn{1}{c|}{worst} \T\B \\
\hhline{|=|=|=|=|=|}
$6$  & $3.66\cdot10^{-3}$ & $3.68\cdot10^{-3}$ & $1.35\cdot10^{-3}$ & $6.17\cdot10^{-3}$ \T\B \\
\hline
$12$ & $2.92\cdot10^{-3}$ & $3.07\cdot10^{-3}$ & $8.95\cdot10^{-4}$ & $9.49\cdot10^{-4}$ \T\B \\
\hline
$24$ & $5.18\cdot10^{-3}$ & $9.53\cdot10^{-3}$ & $9.18\cdot10^{-3}$ & $2.22\cdot10^{-2}$ \B \\
\hline
\end{tabular}
\\[6pt]
\caption{Study~3: open-loop test-set reconstruction accuracy of the aerodynamic loads as a function of the
network depth, with the six-signal input: combined NRMSE, used here to rank configurations
within the sweep, for the architectures of \cref{tab:depth_arch} at the two latent dimensions.
Each configuration is trained twice from two different random initialisations of the weights,
and both outcomes are reported: the better of the two is the LDNet one would retain, the worse
bounds the spread observed over these two runs.}
\label{tab:depth_loads}
\end{table}

\subsubsection{Field reconstruction (Study 4)}
\label{subsubsec:field_recon}

Study~4 evaluates $\mathcal{NN}_\mathrm{rec}$ of \labelcref{eq:reduced_system} at query
coordinates away from the airfoil surface, where it returns the distributed field
$\widetilde{\mathbf{y}}=(v_x,v_y)$ instead of the sectional loads. This capability is not used
by the controller, which queries the single point $\mathbf{x}_\mathrm{query}$; we report it
because the meshless formulation provides it at no additional training cost, and because it
shows where the reconstruction degrades. \Cref{fig:cc050_umag} shows the reconstruction over a
full trajectory on a combined gust\,+\,flap case with gust peak
$W_0=\SI{32.1}{\meter\per\second}$, gust duration $T_g=\SI{1.12}{\second}$ and peak flap
deflection $\SI{5.27}{\degree}$: the LDNet reproduces the velocity field closely at the
quiescent and recovery instants, and least sharply near the airfoil and in the wake at the
peak flap deflection.

\begin{figure}[t]
\centering
\includegraphics[width=\textwidth]{Thesis_article_format/Chapters/Chapter 3/Images/fig_Cc050_recon_umag.png}
\caption{Study~4: CFD (top) and LDNet (bottom) velocity-magnitude field on a
validation trajectory not trained on, with combined gust\,+\,flap, with gust peak
$W_0=\SI{32.1}{\meter\per\second}$, gust duration $T_g=\SI{1.12}{\second}$ and peak flap deflection $\SI{5.27}{\degree}$, at the
quiescent state, the gust peak, the peak flap deflection and recovery. The reconstruction
matches CFD closely except at the peak flap deflection, where it is least sharp near the
airfoil and in the wake.}
\label{fig:cc050_umag}
\end{figure}

\Cref{tab:cost_accuracy} compares the cost of that reconstruction against the full-order CFD
solve that generated its training data, at the same output points (the $11\,075$-node crop
mesh; the CFD solve itself runs on the full $374\,584$-cell production grid of \cref{app:vv}).
This configuration reaches an
NRMSE of $1.10\cdot10^{-2}$ on $v_x$ and $8.32\cdot10^{-3}$ on $v_y$.

\begin{table}[t]
\centering
\begin{tabular}{|c|c|c|}
\hhline{|-|-|-|}
\rowcolor{bluePoli!40}
\multicolumn{1}{|c|}{} & \multicolumn{1}{c|}{CFD (\textit{pimpleFoam})} & \multicolumn{1}{c|}{LDNet reconstruction} \T\B \\
\hhline{|=|=|=|}
cores & $16$ & $1$ \T\B \\
\hline
wall-clock, one trajectory & $\approx\SI{2}{\hour}$ & \SIrange{18.8}{19.5}{\second} \T\B \\
\hline
core-hours, one trajectory & $\approx32$ & $\approx5.3\cdot10^{-3}$ \B \\
\hline
\end{tabular}
\\[6pt]
\caption{Cost of reconstructing one full trajectory ($150$ frames on the $11\,075$-node crop mesh)
with the full-order CFD solve versus the field LDNet of \cref{subsubsec:field_recon}. The CFD cost is the
replay that generated the training/validation/test trajectories themselves, not a
separate benchmark; the LDNet cost is the mean of $10$ repeated forward passes (after a
warm-up call) on a single CPU core.}
\label{tab:cost_accuracy}
\end{table}

\subsubsection{Retained LDNet for control}
\label{subsubsec:rollout_results}

The retained LDNet combines the results of the studies above: the six-signal input of
\cref{fig:input_study} (Study~1) and $d_s=10$, the most accurate latent dimension of
\cref{subsubsec:latent_dim} (Study~2). Depth is the one hyperparameter they do not settle by themselves:
\cref{tab:depth_loads} of Study~3 found $L=12$ the most accurate configuration of the entire study, but in
the open-loop, teacher-forced measure of \cref{subsubsec:depth}. Accuracy in that regime is
not what the controller depends on, since the retained LDNet is not run teacher-forced but in
autonomous rollout, fed only its own predictions once deployed inside the control loop.
Therefore, we compared $L=6$ and $L=12$ directly on the rollout loss
$\mathcal{L}_\mathrm{rollout}$ of \eqref{eq:rollout_loss}, the one on which the retained
LDNet is actually trained and judged, with a single training run per depth rather than
the multiple seeds of Study~3.

Two values of $\mathcal{L}_\mathrm{rollout}$ mark this comparison for each depth: the
warm-start value, from the teacher-forced weights evaluated in autonomous rollout before
any rollout-specific adjustment, and the fine-tuned value, after continuing training on
$\mathcal{L}_\mathrm{rollout}$ itself. At the warm start $L=12$ is already worse than
$L=6$ ($\mathcal{L}_\mathrm{rollout}=0.0864$ against $0.0373$), and fine-tuning widens
rather than closes the gap: $L=6$ reaches $1.79\cdot10^{-3}$, while $L=12$ stalls at
$0.0400$, more than twenty times higher, only after its BFGS optimiser aborts a line
search and is forced to restart. Both values are best-validation checkpoints, not the
end point of a converged run (\cref{fig:loss}(b) details this for $L=6$). Unlike in the open-loop regime, where added depth is a free improvement
in both accuracy and repeatability, in this single comparison it works against the training
objective: a plausible explanation is that a more expressive one-step map is harder to keep
well-behaved once composed with itself for hundreds of steps, though a single seed per depth
does not rule out run-to-run scatter as a contributor. The retained LDNet therefore uses $L=6$: found to be the more stable configuration under autonomous rollout in this comparison, and the one that runs inside the
control loop of the GLA studies of \cref{subsec:mpc_results}.

At $L=6$, the dynamics network $\mathcal{NN}_\mathrm{dyn}$ has two hidden layers of $7$ neurons and
the reconstruction network $\mathcal{NN}_\mathrm{rec}$ four hidden layers of $24$ neurons, all with
$\tanh$ activation and the tail-expanding nonlinear output layer ($\beta=0.05$); the integration
and reference time steps are $\Delta t=\Delta t_\mathrm{ref}=\SI{2e-3}{\second}$. Training uses the
two-stage teacher-forced schedule of \cref{subsubsec:training_algorithm} ($400$ Adam epochs at
learning rate $10^{-2}$ followed by BFGS) with the weight regularisation, the damped latent
update~\eqref{eq:latent_dyn} and the validation early stopping described above.
\Cref{fig:loss}(a) shows the corresponding training history. The
validation curve flattens over the last decade of epochs, reaching
$\mathcal{L}_\mathrm{val}\approx1.3\cdot10^{-4}$ at the retained minimum near epoch $4.2\cdot10^{3}$,
while the training loss keeps decreasing over the same stretch, reaching
$\mathcal{L}_\mathrm{train}\approx3.2\cdot10^{-5}$, about a factor of $4$ below the validation
value at that point: the two track each other closely up to the Adam-to-BFGS
transition, then separate once the Tikhonov term has been driven small, the
onset of overfitting that the validation-patience rule is designed to
catch. The retained weights are those at the validation minimum rather than
at the end of the run, so the LDNet used downstream is the one taken just
before that separation becomes established, not the more overfit one the
unstopped run would reach. Unlike this teacher-forced
pretraining, which in \cref{fig:loss}(a) runs to epoch $4.4\cdot10^{3}$ once the validation-patience
rule above fires past its minimum near epoch $4.2\cdot10^{3}$, the fine-tuning stage that follows keeps the same architecture and inputs,
but skips the Adam warm-up and runs a much shorter schedule: pure BFGS directly
on $\mathcal{L}_\mathrm{rollout}$, saving the weights whenever the rollout loss improves.
\Cref{fig:loss}(b) reports that second stage. Warm-started from the teacher-forced weights, the
rollout loss falls from $3.7\cdot10^{-2}$ to a best validation value of
$1.79\cdot10^{-3}$ at iteration $290$; past that point the validation loss drifts upward again,
reaching $4.5\cdot10^{-3}$ by iteration $500$, while the training loss keeps decreasing
monotonically to $7.0\cdot10^{-4}$. The run is stopped once the budget of $500$ BFGS
iterations is spent rather than at a plateau. The weights retained are therefore
those of iteration $290$.

This divergence between a monotonically falling training loss and a validation loss that
deteriorates past a minimum is an ordinary generalisation gap, sharpened by a mechanism specific
to the autonomous-rollout objective: because each of the $N_\mathrm{rollout}=800$ steps
feeds the next one its own error, the map from weights to rollout loss amplifies small parameter
changes far more violently than the one-step teacher-forced map does. Past the minimum,
this sensitivity gives the optimiser room to keep lowering the training loss by fitting
weight settings that cancel the accumulated error along the specific training trajectories
rather than by improving the underlying one-step dynamics, a fit that has no reason to
transfer to the held-out trajectories of the validation set and that becomes easier to find
the longer training proceeds --- the error-accumulation mechanism identified for autoregressive
neural surrogates~\cite{sanchezgonzalez2020learning,brandstetter2022message}. Retaining the best-validation checkpoint rather than the
final weights is the standard mitigation in that literature, and is what is done here.

\begin{figure}[t]
\centering
\includegraphics[width=0.85\textwidth]{Thesis_article_format/Chapters/Chapter2/Images/fig_ch2_loss.png}
\caption{Retained LDNet ($d_s=10$, six-signal input, $L=6$): the two stages of the training of
the model that runs inside the control loop of \cref{subsec:mpc_results}.
(a) Teacher-forced pretraining: training loss and
validation loss versus epoch. The dashed line marks the transition from the Adam
warm-up to the BFGS phase. Training continues past the minimum of the validation loss until
the plateau criterion fires, and the weights at that minimum are the ones kept. (b) Closed-loop
rollout fine-tuning, warm-started from the weights of (a) and optimised with pure BFGS on
$\mathcal{L}_\mathrm{rollout}$, evaluated every $10$ iterations and shown over the whole run.
The validation loss reaches its minimum at iteration $290$.}
\label{fig:loss}
\end{figure}

% ---- Prepared placeholders (to be written when the field-reconstruction study is ready) ----
% \subsubsection{Latent leak $\lambda$}                       % stability-vs-accuracy trade-off (train_lambdas.sh)

\subsection{GLA over the discrete-gust envelope}
\label{subsec:mpc_results}

The controller is evaluated over a parametric grid of 18 gust cases inspired by the
CS-25.341 discrete gust: gust amplitude $W_0 \in \{10,\,20,\,30\}\,\si{\meter\per\second}$
and gust duration $T_g \in \{0.3,\,0.4,\,0.5,\,0.7,\,1.0,\,1.2\}\,\si{\second}$, swept directly
rather than derived from the $H$-dependent design gust velocity $U_{ds}$ that CS-25.341
prescribes, restricted to $H\le\SI{157}{\foot}$ of its $H\le\SI{350}{\foot}$ range and,
through the sign of \labelcref{eq:gust}, to upward gusts alone.
Each case is located in the aerodynamically relevant parameter space by two quantities derived
from the gust parameters of \cref{eq:gust}: the gust gradient distance $H = U_\infty T_g/2$, the
length with which CS-25.341 parametrises the discrete-gust envelope and is therefore reported in
feet, and the reduced frequency $k = \pi\,c_{\mathrm{ref}}/(U_\infty T_g)$.

The primary performance metric is the fractional reduction of the peak
lift-coefficient excursion with respect to the open-loop reference:
\begin{equation}
    \mathrm{CLred} = \frac{\max_t |C_L^{\mathrm{OL}}(t) - C_{L,\mathrm{trim}}|
                         - \max_t |C_L^{\mathrm{CL}}(t) - C_{L,\mathrm{trim}}|}
                         {\max_t |C_L^{\mathrm{OL}}(t) - C_{L,\mathrm{trim}}|}
                    \times 100\%,
\label{eq:clred}
\end{equation}
where the maximum is taken over the gust window $t \in [0,\, T_g + \SI{0.5}{\second}]$ and
$C_{L,\mathrm{trim}}$ is the target the controller drives the lift towards, the same
reference whose deviations the MPC cost~\eqref{eq:mpc_cost} penalises: CLred is therefore
the fraction of the open-loop excursion about that target that the controller removes.
A secondary stability check excludes any tuning that causes any of
$\dot\alpha$, $\ddot\alpha$, or $\ddot h$ to exceed three times the
corresponding open-loop peak, a conventional margin rather than a certification limit.
For each gust case the effort weight
of the cost~\eqref{eq:mpc_cost} is swept over
$R \in \{10^{-4},\,3\cdot10^{-4},\,10^{-3},\,3\cdot10^{-3},\,10^{-2}\}$,
five half-decade steps, and the selected value, written $R^{\star}$, is the
one maximising CLred among those that satisfy the check. The weight is
therefore tuned case by case rather than held fixed over the envelope, and
$R^{\star}$ is reported for every case alongside CLred.
The pitch ratio, defined as the ratio of the closed-loop to open-loop peak
pitch excursion, is monitored alongside CLred to make explicit any structural
load trade-off introduced by the flap. In this nominal evaluation the sensing
chain is noise-free, so the preview vector coincides with the sampled true
gust profile~\eqref{eq:gust}.
The numerical values of the structural, aerodynamic and actuator parameters shared by every
run of this paragraph are collected in \cref{tab:results_setup} of \cref{app:fom}; the symbols
are those of the structural model~\eqref{eq:eom}--\eqref{eq:matrices}. The structural and flap
values are adapted from the typical-section aeroelastic benchmark of Hodges and
Pierce~\cite{hodges2011aeroelasticity}.

\subsubsection{Performance over the discrete-gust envelope}

\Cref{fig:mpc_traces} shows the nominal closed-loop response at the two gust cases this work
follows throughout: Test~1 ($W_0 = \SI{20}{\meter\per\second}$, $T_g = \SI{0.7}{\second}$, the gust
of \cref{fig:cfd}) and Test~2 ($W_0 = \SI{10}{\meter\per\second}$, $T_g = \SI{0.4}{\second}$, the
short gust of \cref{subsec:noise_results,subsec:systematic_results}). The preview lets the flap
begin acting while the gust is still building up, first countering the initial lift dip at gust
onset and then reversing against the main excursion; the closed-loop lift stays confined to a
narrow band around trim, with at most $\SI{8.4}{\degree}$ and $\SI{4.2}{\degree}$ of deflection
and pitch excursions comparable to or moderately larger than open loop (ratios $0.90$ and
$1.15$); the peak-excursion reduction is reported in the caption of \cref{fig:mpc_traces}.


\Cref{tab:mpc_cost_accuracy} times both ROM and FOM closed loops on Test~1: the full-order loop
is roughly $46\times$ slower in wall-clock time and $\approx700\times$ in core-hours, and against
this gap the accuracy cost is modest, an $11.2$-point CLred difference between the ROM and FOM
controlled systems, both driven by the same LDNet internal model, on Test~1 and $14.5$ points
on Test~2, the only two cases run full-order. The ROM
entries are an interpreted research prototype on a single core, not a deployment-oriented
implementation: at $\Delta t = \SI{2}{\milli\second}$ one simulated second requires $500$
decisions, so the measured \SI{80.2}{\second} per simulated second is about \SI{160}{\milli\second}
per decision, roughly $80\times$ the \SI{2}{\milli\second} the step allows, spent on the exhaustive
$161$-candidate search rather than
on the LDNet itself. This work therefore establishes the cost reduction against the full-order
model, not embedded real-time feasibility, which would need a compiled implementation and a
cheaper search.


\Cref{tab:cs25_grid} reports twelve of the eighteen cases, at $T_g\in\{0.3,0.4,0.7,1.2\}\,\si{\second}$
for each amplitude, the four durations that span the gradient range; the full eighteen-case
grid, including $T_g\in\{0.5,1.0\}\,\si{\second}$, is archived in the repository. At
$T_g = \SI{0.7}{\second}$, the middle of the gradient range shown in \cref{tab:cs25_grid}, the
controller removes more than $91\%$ of the peak excursion at every amplitude, the highest and
most uniform reduction on the grid. Performance degrades toward both
ends of the envelope, for different reasons. At the sharp--severe corner
($W_0 = \SI{30}{\meter\per\second}$, $T_g = \SI{0.3}{\second}$) the reduction is limited to $39.4\%$ and is
only weakly sensitive to the effort weight (below $5$ percentage points of
variation over two decades of $R$, with no stability violation): the peak
gust slope $\pi W_0/T_g$ corresponds to an equivalent-incidence rate
$\dot{W}/U_\infty \approx \SI{225}{\degree\per\second}$, comparable to the actuator rate
limit itself, so the flap---aerodynamically less effective than
incidence---cannot track the gust ramp regardless of the tuning. At the
long-gust corner the stability check becomes the active constraint: the
large sustained deflection required to cancel a quasi-steady gust drives the
pitch degree of freedom, whose open-loop peaks are small for slow gusts;
aggressive tunings amplify the pitch excursion (ratios up to $1.9$) and
violate the $\dot\alpha$ bound, so the selected weight rises to
$R^{\star} = 10^{-2}$ and the reduction settles at $41.7\%$
($W_0 = \SI{30}{\meter\per\second}$) and $58.2\%$ ($W_0 = \SI{20}{\meter\per\second}$) at $T_g = \SI{1.2}{\second}$.

\begin{figure}[H]
    \centering
    \includegraphics[width=\textwidth]{Thesis_article_format/Chapters/Chapter 3/Images/fig_ch3_trace_tests.png}
    \caption{Noise-free closed-loop response at the two gust cases of the CS-25
    grid of \cref{tab:cs25_grid} retained throughout this work, Test~1
    ($W_0 = \SI{20}{\meter\per\second}$, $T_g = \SI{0.7}{\second}$, left) and Test~2 ($W_0 = \SI{10}{\meter\per\second}$,
    $T_g = \SI{0.4}{\second}$, right): gust profile (top), lift coefficient (middle) and
    flap deflection (bottom), over the evaluation window
    $t \in [0,\,T_g + \SI{0.5}{\second}]$. Solid lines are the reduced-order
    system (ROM): open loop (red), closed loop (light blue), and the
    commanded flap (orange), against the trim value $C_{L,\mathrm{trim}}$
    (dotted). Dashed lines repeat the closed loop on the full-order model
    (FOM): the controller still predicts with the LDNet surrogate, but is now
    fed the true structural state of the co-simulation at every coupling
    window and re-plans against it, so its commanded deflection genuinely
    differs from the reduced-order one rather than replaying it. The
    full-order closed loop reduces the peak excursion by $80.7\%$ (Test~1) and $74.6\%$
    (Test~2), against $91.9\%$ and $89.1\%$ for the reduced-order loop, each case at its
    own $R^\star$ of \cref{tab:cs25_grid}. Both reductions are computed against their own
    plant's open-loop peak; the full-order open-loop trace is not plotted; reading the
    full-order reduction off the figure by comparing the dashed closed loop against the
    solid red (reduced-order) open loop underestimates it, because the ROM itself
    underestimates the open-loop peak by $28\%$ (Test~1) and $44\%$ (Test~2).}
    \label{fig:mpc_traces}
\end{figure}

Across the
whole grid the deflection limit is never reached ($|\delta|_{\max} \le \SI{12.2}{\degree} <
\delta_{\max}$; the rate limit was not separately audited) and no selected tuning violates the stability check; the
pitch ratio spans $0.57$--$1.47$. This peak magnitude exceeds the \SI{8.5}{\degree}
deflection bound of family Cc of \cref{tab:dataset_families}, the only training family
where gust and flap are correlated as they are in closed loop, so the LDNet is queried
somewhat beyond that family's deflection range at the more demanding cases of the grid,
a point the closed-loop results do not otherwise flag.

\subsection{Robustness to lidar measurement noise}
\label{subsec:noise_results}

The stochastic robustness of the pipeline is assessed with both noise
models of \cref{subsec:mpc_arch}: the constant-level model drives the
amplitude sweep of this paragraph, while the
range-proportional lidar model~\eqref{eq:range_prop_noise} enters the single-case check below. Both studies of this paragraph and of \cref{subsec:systematic_results} are run on Test~2.

\begin{table}[t]
\centering
\begin{tabular}{|c|c|c|}
\hhline{|-|-|-|}
\rowcolor{bluePoli!40}
 & ROM & FOM \T\B \\
\hhline{|=|=|=|}
cores & $1$ & $16$ \T\B \\
\hline
wall-clock per simulated second & \SI{80.2}{\second} & \SI{3.68e3}{\second} \T\B \\
\hline
core-hours per simulated second & $2.2\cdot10^{-2}$ & $16.3$ \T\B \\
\hline
CLred & $91.9\%$ & $80.7\%$ \B \\
\hline
\end{tabular}
\\[6pt]
\caption{Cost of one closed-loop second at Test~1 ($W_0=\SI{20}{\meter\per\second}$,
$T_g=\SI{0.7}{\second}$), reduced-order (ROM) versus
full-order (FOM) controller. ROM cost is the wall-clock time of the controller-plus-plant
loop of \cref{subsec:mpc_arch} (oracle preview, $N=8$, $161$-point flap search), timed on a
single CPU core. FOM cost is measured on a short freshly-timed segment of the co-simulation
of \cref{fig:mpc_traces} (the controller re-plans once per $\SI{3.5}{\milli\second}$
coupling window of \cref{subsubsec:temporal}, coarser than its own $\SI{2}{\milli\second}$
re-planning step, $16$ cores). The wall-clock ratio is $\approx46\times$; the
core-hour ratio, accounting for the FOM's $16$ cores against the ROM's one, is
$\approx7.3\cdot10^{2}\times$.}
\label{tab:mpc_cost_accuracy}
\end{table}

\begin{table}[H]
    \centering
    \begin{tabular}{|r|r|r|r|r|r|r|r|}
    \hline
    \rowcolor{bluePoli!40}
    $W_0$ [\si{\meter\per\second}] & $T_g$ [\si{\second}] & $H$ [\si{\foot}] & $k$ & CLred [\si{\percent}] & $R^{\star}$ &
    $|\delta|_{\max}$ [\si{\degree}] & pitch ratio \T\B \\
    \hline\hline
    10 & 0.30 &  39 & 0.131 & 81.2 & $10^{-3}$          & 4.6  & 1.47 \T\B \\
    \hline
    10 & 0.40 &  52 & 0.098 & 89.1 & $3\times 10^{-4}$  & 4.2  & 1.15 \T\B \\
    \hline
    10 & 0.70 &  92 & 0.056 & 93.5 & $3\times 10^{-4}$  & 3.3  & 0.96 \T\B \\
    \hline
    10 & 1.20 & 157 & 0.033 & 93.4 & $10^{-3}$          & 2.5  & 0.96 \T\B \\
    \hline\hline
    20 & 0.30 &  39 & 0.131 & 86.0 & $3\times 10^{-4}$  & 8.8  & 0.57 \T\B \\
    \hline
    20 & 0.40 &  52 & 0.098 & 87.8 & $10^{-4}$          & 8.4  & 1.08 \T\B \\
    \hline
    20 & 0.70 &  92 & 0.056 & 91.9 & $10^{-4}$          & 8.4  & 0.90 \T\B \\
    \hline
    20 & 1.20 & 157 & 0.033 & 58.2 & $10^{-2}$          & 5.4  & 1.14 \T\B \\
    \hline\hline
    30 & 0.30 &  39 & 0.131 & 39.4 & $10^{-4}$          & 9.1  & 0.87 \T\B \\
    \hline
    30 & 0.40 &  52 & 0.098 & 80.5 & $3\times 10^{-4}$  & 7.9  & 1.02 \T\B \\
    \hline
    30 & 0.70 &  92 & 0.056 & 91.8 & $10^{-4}$          & 10.5 & 0.75 \T\B \\
    \hline
    30 & 1.20 & 157 & 0.033 & 41.7 & $10^{-2}$          & 7.7  & 1.27 \B \\
    \hline
    \end{tabular}
    \\[6pt]
    \caption{Noise-free results at twelve of the eighteen CS-25.341 gust cases,
    $T_g\in\{0.3,0.4,0.7,1.2\}\,\si{\second}$ at each amplitude: gust
    gradient distance $H = U_\infty T_g/2$ and reduced frequency
    $k = \pi\,c_{\mathrm{ref}}/(U_\infty T_g)$, peak lift-excursion reduction of
    \cref{eq:clred}, effort weight $R^{\star}$ selected per case, peak flap
    deflection and pitch ratio. No case violates the stability check.}
    \label{tab:cs25_grid}
\end{table}


The
sweep applies
the measurement model~\eqref{eq:meas_model} with the constant per-sample
level $\sigma_j = \sigma$: at every control instant $t_n$ each of the
$J_{\max}$ raw samples is corrupted by its own error draw,
\begin{equation*}
y_j(t_n) = W\bigl(t_n + j\,\Delta t\bigr) + \varepsilon_{nj}, \qquad
\varepsilon_{nj} \sim \mathcal{N}\bigl(0,\,\sigma^2\bigr), \qquad
\frac{\sigma}{W_0} \in \{0, 1, 2, 5, 10, 20\}\%,
\end{equation*}
with the errors drawn independently across nodes and control instants:
white noise of up to \SI{2}{\meter\per\second} per shot. The noisy samples enter the
fusion sums of \cref{eq:invvar} exactly as drawn, and the aeroelastic
system keeps advancing with the true gust $W(t)$: the noise reaches the
control loop only through the preview vector $\widehat{W}_k$ in the
horizon cost of \cref{eq:mpc_cost}. Since the outcome depends on the particular random
sequence drawn, we simulate each level six times with independent noise
realisations, and take the statistics over the six repetitions.

\Cref{fig:noise_white} reports the resulting reduction together with the
error actually delivered to the controller after fusion,
$\sigma_{\mathrm{del}} = \mathrm{std}\bigl(\widehat{W}_1(t) - W(t+\Delta t)\bigr)$,
the standard deviation of the residual on the first preview node over the
control steps of the evaluation window, averaged over the realisations.
Since each gust
value is measured $J_{\max} = 50$ times with independent errors before
reaching the wing, the inverse-variance average of \cref{eq:invvar}
shrinks the per-sample error by roughly $\sqrt{J_{\max}} \approx 7$: at
$\sigma = \SI{0.2}{\meter\per\second}$ the delivered error is $\SI{0.02}{\meter\per\second}$, about one order of
magnitude below the raw level. This reduction rests on the measurement model of
\cref{subsec:controller}: the gust convects rigidly, so the $J_{\max}$ samples of a given
gust value are repeated measurements of one quantity with independently drawn errors. Errors
correlated along the range, which airborne lidar exhibits through coherence loss and finite
probe volume, would not average down in the same proportion, so the tolerance reported here
is a property of the loop under this measurement model rather than a characterisation of
lidar hardware. The noise that actually reaches the horizon
cost~\eqref{eq:mpc_cost} is therefore far smaller than the nominal
$\sigma$, and as a result the mean reduction remains within about one
percentage point of the noise-free value up to $\sigma = 10\%\,W_0 =
\SI{1}{\meter\per\second}$---the lower end of the per-shot error of airborne Doppler hardware
discussed in \cref{subsec:mpc_arch}. Only at $\sigma = 20\%\,W_0$ does the
spread across realisations open appreciably (worst realisation $82.1\%$,
none flagged).

\begin{figure}[t]
    \centering
    \includegraphics[width=0.7\textwidth]{Thesis_article_format/Chapters/Chapter 3/Images/fig_ch3_noise_white_W10.png}
    \caption{Peak lift-coefficient excursion reduction under white measurement noise on the raw
    lidar samples, on Test~2 ($W_0 = \SI{10}{\meter\per\second}$, $T_g = \SI{0.4}{\second}$): mean
    and min--max band over 6 noise realisations, with the delivered
    post-fusion error $\sigma_{\mathrm{del}}$ on the second axis row (a ringed
    point, unused here, would mark the stability flag raised in the indicated
    fraction of realisations). The mean reduction stays within about one
    percentage point of the noise-free $89.1\%$ up to $\sigma = 10\%\,W_0$,
    and only at $20\%\,W_0$ does the spread across realisations open
    appreciably, with no stability flag raised at any level.}
    \label{fig:noise_white}
\end{figure}

A further check verifies that this tolerance also holds under the
range-proportional lidar model~\eqref{eq:range_prop_noise}, whose raw
line-of-sight error (\SIrange{1}{3}{\meter\per\second}) is largest relative to the gust on Test~2 of
any case in the study: with the one configuration already tuned for the
noise-free case ($R = R^\star = 3\times 10^{-4}$, left unchanged for every
noise and lidar scenario), we run Test~2 under white noise at $\sigma =
2\%\,W_0$ and under this lidar model, and in both cases the mean reduction
stays within $3.5$ percentage points of the noise-free value, with no
stability flag in any realisation (noise-free/white/lidar CLred is
$89.1\%/89.1\%/85.6\%$).

\subsection{Systematic sensing errors}
\label{subsec:systematic_results}

Zero-mean noise is, however, the error the fusion layer is built to remove;
systematic errors pass through it unaffected. On the sensing side they are
modelled by corrupting the raw samples of \cref{eq:meas_model} as
\begin{equation*}
y_j(t_n) = W\bigl(t_n + j\,\Delta t + \tau\bigr) + b + \varepsilon_{nj},
\end{equation*}
with one error active at a time: a calibration bias $b = \beta\,W_0$,
identical for every sample, and a time shift $\tau$ of the measured
profile, positive when the gust appears in the preview earlier
than it arrives. Since the fusion
\eqref{eq:invvar}--\eqref{eq:tikhonov} is linear in the samples, averaging
leaves both terms untouched: the delivered preview is
$\widehat{W}_k \approx W(t + k\Delta t + \tau) + b$, an error that
enters every term of the horizon cost~\eqref{eq:mpc_cost} coherently
instead of averaging down.

\Cref{fig:noise_bias_timing} sweeps the two errors, again on Test~2. In
both panels the curve is a deterministic sweep with the systematic error
alone ($\varepsilon_{nj} = 0$); at $b = \pm 5\%\,W_0$ and
$\tau = \pm\SI{10}{\milli\second}$ the white noise of the previous paragraph
($\sigma = 2\%\,W_0$, six realisations) is superimposed on the systematic
error to probe the combination of the two error classes. On the tolerated
side the noisy means reproduce the deterministic values within a fraction
of a percentage point with no flag, so the two classes do not interact.
The calibration axis is strongly asymmetric. A calibration bias that
over-reads the gust is nearly free at $+2\%\,W_0$ (CLred $90.2\%$) but already degrades
substantially by $+5\%\,W_0$ ($53.2\%$, no flag), and collapses at $+10\%\,W_0$
(flagged, deterministic value below zero), while under-reading is nearly
free down to $-10\%\,W_0$ (CLred $90.0\%$, no flag): an over-read gust keeps the
flap deployed against a disturbance that never materialises, driving the
flap excursion and the pitch peak up as the bias grows, until the
self-induced transient destabilises the loop. For the systematic time
shift of the preview, the tolerated range covers the entire sweep run,
$\tau \in$ \SIrange{-20}{20}{\milli\second}, with no flag raised at either
end, so the timing axis is tolerated symmetrically; the two directions are
nonetheless not symmetric in value: a delayed preview does not merely
avoid degrading the reduction, it is mildly beneficial, $+2.2$ percentage
points at $\SI{-20}{\milli\second}$ (CLred $91.3\%$, peaking at $91.6\%$ at
$\SI{-10}{\milli\second}$), while an anticipated preview of
$\SI{+20}{\milli\second}$ costs $3.9$ points (CLred $85.2\%$); the receding
horizon re-plans from the measured state at every step, which keeps both
directions well short of the instability seen on the calibration axis over
the same sweep range.



\begin{figure}[t]
    \centering
    \includegraphics[width=1\textwidth]{Thesis_article_format/Chapters/Chapter 3/Images/fig_ch3_bias_timing_W10.png}
    \caption{Robustness to the two critical systematic errors, on Test~2
    ($W_0 = \SI{10}{\meter\per\second}$, $T_g = \SI{0.4}{\second}$): calibration bias of the wind
    measurement (left) and systematic preview time shift (right).
    Deterministic noise-free sweeps with spot checks at $\sigma = 2\%\,W_0$
    (6 noise realisations, min--max bars); ringed points raise the
    stability flag. The calibration axis is strongly asymmetric:
    over-reading the gust degrades the reduction well before $+10\%\,W_0$,
    where it collapses and raises the stability flag, whereas
    under-reading it is nearly free down to $-10\%\,W_0$. The timing axis
    is tolerated symmetrically over the sweep range shown, with no flag
    raised at either end.}
    \label{fig:noise_bias_timing}
\end{figure}

Both errors act on the preview, and therefore on the command alone: the internal model and
the controlled system share the same LDNet throughout this paragraph. The sensitivity to the
internal model itself is the subject of \cref{app:modelcomp}, which complements the
full-order closed loop of \cref{fig:mpc_traces}.

\clearpage
